\documentclass[lang=cn,12pt,chinesefont=founder]{elegantbook}
\usepackage{graphicx}
\usepackage{unicode-math}
\usepackage{framed} %段落加边框
\usepackage{makecell} %表格内强制换行
\usepackage{titletoc} % 章节目录
\usepackage{tabularx}
\usepackage{array}
\setmathfont{XITS Math}


\begin{document}
\chapter{LG 15Z980黑苹果安装}

\section{Bios设置}

参考自17z990配置，部分选项没有 \link{https://github.com/lehrian/LG-Gram-17z990-Hackintosh}

Main：
    \begin{itemize}
        \item Boot Features
            \begin{itemize}
            \item Quick Boot => Disabled
            \end{itemize}
    \end{itemize}

Advanced：

1. SW Guard Extensions (SGX) => Disabled

2. Intel Advanced Menu
            \begin{itemize}
            \item CPU Configuration
                \begin{itemize}
                \item Software Guard Extensions (SGX) => Disabled
                \end{itemize}
            \item Power  Performance
                \begin{itemize}
                \item CPU - Power Management Control
                \begin{itemize}
                  \item CPU Lock Configuration
                  \begin{itemize}
                    \item CFG Lock => Disabled
                    \item Overclocking Lock => Disabled
                  \end{itemize}
                \end{itemize}
              \end{itemize}
            \item System Agent (SA) Configuration
                \begin{itemize}
                \item Graphics Configuration
                \begin{itemize}
                  \item DVMT Pre-Allocated => 64M
                \end{itemize}
              \end{itemize}
            \item PCH-IO Configuration
                \begin{itemize}
                \item Security Configuration
                \begin{itemize}
                  \item Force unlock on all GPIO pads => Enabled (Thanks 1OldSWguy)        
                \end{itemize}
              \end{itemize}
            \item Platform Settings
                \begin{itemize}
                \item System Time and Alarm Source => Legacy RTC
              \end{itemize}
            \item Thunderbolt Configuration
                \begin{itemize}
                  \item Thunderbolt Boot Support => Pre-Boot ACL
                  \item Security Level => No Security
                \end{itemize}
            \end{itemize}

\section{启动U盘制作}

\begin{enumerate}
  \item 下载EFI \link{https://github.com/a31user/LG-15Z980-Monterey-EFI}
  \item 下载系统镜像 \link{https://macoshome.com/macos}
  \item 下载balenaEtcher \link{https://www.balena.io/etcher}
  \item 制作macOS启动盘 \link{https://zhuanlan.zhihu.com/p/521391783}
   \begin{itemize}
     \item 插入要刻录为macOS启动盘的优盘，打开balenaEtcher。
     \item 点击Flash from file，选择macOS的安装dmg文件，选择U盘，点击Flash
     \item 等待写入并验证完成，如果弹出需要格式化，一定要点击取消，因为Windows无法识别macOS安装U盘。
     \item 打开分区工具DiskGenius，给U盘的EFI分区分配盘符
     \item 将下载的EFI文件夹复制到EFI分区中
   \end{itemize}
   \item 下载恢复镜像 \link{http://imacos.top/2022/12/01/macos-base-system-recovery-2/}
   \item 打开分区工具 DiskGenius，给 U 盘的未使用空间新建分区，fat32格式
   \item 将com.apple.recovery.boot文件夹复制到新建分区中（包含.contentDetails）
\end{enumerate}


\section{系统装好之后的工作}

\subsection{EFI分区挂载}

使用OpenCore configurator.app

\begin{enumerate}
  \item 打开OpenCore configurator，工具-挂载EFI。挂载U盘EFI分区
  \item 把U盘里面的 EFI 文件夹复制到电脑桌面
  \item 挂载系统盘EFI分区
  \item 把桌面EFI 文件夹复制到系统盘EFI分区
\end{enumerate}

\subsection{开机跑代码怎么关掉}

修改Config.plist

Windows下用notepad3打开，搜索-v，删掉。

不能用OpenCore configurator修改，版本不对应。

\subsection{将5秒等待改短}

修改Config.plist

Windows下用notepad3打开，搜索<Key>Timeout<Key>，将5改成2。

\subsection{关掉大版本更新 \link{https://sysin.org/blog/disable-macos-monterey-update/}}

编辑 hosts 文件，添加如下内容

手动编辑：打开终端，执行命令 sudo vi /etc/hosts，添加以下条目：

\begin{lstlisting}
    ## Mac Software Update (sysin)
    127.0.0.1 swdist.apple.com
    127.0.0.1 swscan.apple.com
    127.0.0.1 swcdn.apple.com
    127.0.0.1 gdmf.apple.com
    127.0.0.1 mesu.apple.com
    127.0.0.1 xp.apple.com
\end{lstlisting}

推荐使用 SwitchHosts（免费软件）修改hosts。注意：SwitchHosts 最后要空一行。

例外：如果使用了系统代理或者虚拟专用网，将绕过本地 DNS 设置，仅仅设置 hosts 仍然会检测到更新。

\begin{itemize}
  \item 清除系统更新标记
\end{itemize}

如果已经检测到更新，可以使用如下方法临时去除更新通知标记。

打开 “终端”，执行如下命令：

\begin{lstlisting}
    defaults write com.apple.systempreferences AttentionPrefBundleIDs 0
    Killall Dock
\end{lstlisting}

\subsection{取消密码位数限制}

\begin{enumerate}
  \item 终端中输入\verb|pwpolicy -clearaccountpolicies|
  \item 输入开机密码，然后回车。
  \item 提示 Clearing global account policies 代表成功，否则请检查是否输入正确。
  \item 打开系统偏好设置 > 用户与群组
  \item 更改密码
\end{enumerate}



\begin{lstlisting}
    pwpolicy -clearaccountpolicies
\end{lstlisting}

\section{疑难杂症}
\subsection{不能从U盘启动}
    关闭bios中的secure boot
\subsection{从恢复模式安装}
从此位置下载的EFI只能从恢复模式安装

\begin{lstlisting}
    https://github.com/a31user/LG-15Z980-Monterey-EFI
\end{lstlisting}

\subsection{制作的恢复模式U盘oc菜单没有引导项}

\begin{lstlisting}
    http://www.imacosx.cn/7467.html#toc_1
\end{lstlisting}

先制作完整启动U盘，然后新建分区，把恢复文件放进去

\subsection{恢复模式没有wifi图标，不能联网}

intel网卡驱动AirportItlwm分版本，下载对应系统版本的驱动

\begin{lstlisting}
    https://github.com/OpenIntelWireless/itlwm/releases
\end{lstlisting}

\section{安装情况}
\begin{enumerate}
  \item 全盘安装
  
    \begin{itemize}
        \item Monterey 安装至12分钟重启
    \end{itemize}
  \item 恢复安装
    \begin{itemize}
        \item Catalina 正常安装完成，进系统安装安全更新后wifi不能用
        \item Monterey 正常安装完成，wifi正常。无法系统更新至Ventura
        \item Ventura 卡代码，无法到安装界面
    \end{itemize}
\end{enumerate}






\end{document}